% !TeX root = ../main.tex

%------------------------------- 2
\chapter{基本理论与方法介绍}
	
	高能重离子碰撞为我们提供以实验室条件下研究QGP性质的独特机会。这种研究有助于我们深入了解QCD在极端条件下的行为，并揭示了宇宙早期宏观物质状态的性质。
	在研究高能重离子碰撞的喷注淬火时，由原子核PDF带来的初态冷核物质效应总是与末态的热核物质效应耦合在一起，体现在核修正因子$R_{AA}$中。因此了解冷核物质效应为精确的研究热核物质效应提供了基准，可以更好地理解和分离出QGP形成后的动力学信号。

	
%------------------------------- 2.1
\section{微扰QCD理论}
\label{sec:nlojet_cs}

	由于孤立的部分子没有被观测到，因此需要进行QCD因子化，以便将不同标度的物理过程分离开来，并将硬碰撞中的夸克和胶子与观测到的强子联系起来。因子化方法已经成功应用于对撞机和固定靶实验中。特别是，可以通过全局分析从不同的过程和实验中一致地提取出普适的部子分布函数。
	QCD理论的核心任务是在微扰理论的高阶下，提供各种相关部分子硬散射过程（即涉及到大的转移动量$Q$的过程）的截面和分裂函数的精确计算。过去几年在这个领域取得了巨大的进展。例如，对于纵向动量分布函数，相关的DGLAP演化核现在已经完全计算到了次次领头阶（next-to-next-to-leading order，简称NNLO）~\parencite{Moch:2014sna,Vogt:2004mw,Moch:2004pa}。
	作为强子物理中的基本理论，微扰QCD理论提供了一种定量描述夸克和胶子相互作用的方法，尤其适用于高能散射和高能核物理中的过程。它的核心思想是将夸克和胶子场视为量子场，并利用微扰展开的方法进行计算。

	在~\ref{sec:SM}中已做介绍，根据QCD渐近自由的性质，对于大动量转移或短距离的过程，部分子之间的跑动耦合会减弱，即$\alpha_s(Q^2 \to \infty) \to 0$，此时可以用微扰量子色动力学（Perturbative Quantum Chromodynamics，简称pQCD）来计算。
	例如，QCD喷注是由强子和核碰撞中夸克和胶子的硬碰撞所产生的。这个过程可以通过pQCD很好地描述~\parencite{CMS:2016jip,ATLAS:2017ble,Dokshitzer:1978hw}。
	但是当涉及到强子的产生及其性质时，pQCD理论不再适用，此时是非微扰的动力学机制。
	在计算强子散射截面的pQCD方法中，基于的是夸克部分子模型图景。根据这一图景，任何硬散射过程的截面可以写成夸克和胶子部分子的结构函数、碎裂函数和一个硬散射截面因子的卷积。
	pQCD的因子化定理将能够微扰计算的硬过程和非微扰的软过程分离开，从而实现对两部分的单独计算~\parencite{CTEQ:1993hwr,Collins:1989gx}，
	如图~\ref{fig:perturb}所示。
	在核子-核子碰撞$A+B\to C+X$过程中，高横动量单举强子产生的散射截面为：
\begin{equation}
	\begin{array}{r}
		d \sigma_{A B \rightarrow C X} \approx \int d x_a \int d x_b \int d z_c f_{a / A}\left(x_a, \mu_f\right) \otimes f_{b / B}\left(x_b, \mu_f\right) \\
		\otimes d \sigma_{a b \rightarrow c d}\left(\mu_f, \mu_F, \mu_r\right) \otimes D_{C / c}\left(z_c, \mu_F\right)
	\end{array}
\end{equation}
	其中，$\sigma_{ab\to cX} (x_1,x_2,Q^2)$是微扰部分的散射截面，是短程部分，它可以在强耦合常数$\alpha_s$的特定阶数进行计算。
	部分子分布函数$f_{a/A}(x,Q^2)$表示在核A中发现味道为$a$的部分子且其动量分数为$ x = p_{parton}/p_{nucleus}$的概率。这个量包含长程作用，在因子化框架下跟微扰计算的阶数相对应，它随标度的演化可以用DGLAP方程进行微扰计算，我们可以在任何基本过程中进行测量，比如深度非弹性散射。
	碎裂函数（fragmentation function，简称FF）$D_{c\to h}(z,Q^2)$ 描述出射部分子$c$碎裂为末态强子$h$且其动量分数为$z=p_{hadron}/p_{parton}$的概率~\parencite{deFlorian:2007aj}。
	$z_c$ 为末态强子碎裂动量分数。
	$\mu_f$、$\mu_F$和$\mu_r$分别为因子化的分布函数标度、碎裂函数标度和重整化标度，在理论计算时常取三者相等且与末态领头强子的横动量成正比。
	
	  \hspace*{\fill}
	\begin{figure}[htbp]
		\centering
		\includegraphics[width=3.3in]{perturb}
		\caption{在高能强子-强子碰撞中，双喷注的产生和微扰量子色动力学的共线因子化~\parencite{dEnterria:2009xfs}。}
		\label{fig:perturb}
	\end{figure}

	在微扰QCD理论中，物理量的计算通常采用费曼图的方法。
	\begin{figure}[htbp]
		\centering
		\includegraphics[width=2.0in]{Fenym}
		\caption{领头阶下$qq \to qq$ 过程的费曼图。}
		\label{fig:Fenym}
	\end{figure}
	接下来介绍领头阶的计算，此时 $\sigma_{ab\to cX} (x_1,x_2,Q^2)$ 主要来自 $2\to 2$过程的贡献~\parencite{Owens:1986mp,Ellis:1985er}，以$qq'\to qq'$的过程为例，通常取轻夸克$u, d, s$的质量约为0，费曼图为~\ref{Fenym1}，曼德尔斯通（Mandelstam）变量与末态四动量的关系为：
\begin{equation}
	\begin{aligned}
		& \hat{s}=\left(p_1+p_2\right)^2, \\
		& \hat{t}=\left(p_1-p_3\right)^2, \\
		& \hat{u}=\left(p_1-p_4\right)^2 .
	\end{aligned}
\end{equation}
此过程的微分散射截面为：
\begin{equation}
	\frac{d \sigma}{d \hat{t}}\left(q q^{\prime} \rightarrow q q^{\prime}\right)=\frac{4 \pi \alpha_s^2}{9 \hat{s}^2}\left[\frac{\hat{s}^2+\hat{u}^2}{\hat{t}^2}\right]
\end{equation}
	其它 $2\to 2$的过程结果如下。具体可以参阅文献~\parencite{Owens:1986mp,Ellis:1985er,Babcock:1978gd}，此处不再细致推导。
\begin{equation}
	\begin{aligned}
		\frac{d \sigma}{d \hat{t}}\left(q \bar{q} \rightarrow q^{\prime} \bar{q}^{\prime}\right) & =\frac{4 \pi \alpha_s^2}{9 \hat{s}^2}\left[\frac{\hat{t}^2+\hat{u}^2}{\hat{s}^2}\right], \\
		\frac{d \sigma}{d \hat{t}}(q \bar{q} \rightarrow q \bar{q}) & =\frac{4 \pi \alpha_s^2}{9 \hat{s}^2}\left[\frac{\hat{u}^2+\hat{s}^2}{\hat{t}^2}+\frac{\hat{t}^2+\hat{u}^2}{\hat{s}^2}-\frac{2}{3} \frac{\hat{u}^2}{\hat{s} \hat{t}}\right], \\
		\frac{d \sigma}{d \hat{t}}(q q \rightarrow q q) & =\frac{4 \pi \alpha_s^2}{9 \hat{s}^2}\left[\frac{\hat{u}^2+\hat{s}^2}{\hat{t}^2}+\frac{\hat{t}^2+\hat{s}^2}{\hat{u}^2}-\frac{2}{3} \frac{\hat{s}^2}{\hat{u} \hat{t}}\right], \\
		\frac{d \sigma}{d \hat{t}}(q \bar{q} \rightarrow g g) & =\frac{32 \pi \alpha_s^2}{27 \hat{s}^2}\left[\frac{\hat{u}}{\hat{t}}+\frac{\hat{t}}{\hat{u}}-\frac{9}{4}\left(\frac{\hat{u}^2+\hat{t}^2}{\hat{s}^2}\right)\right], \\
		\frac{d \sigma}{d \hat{t}}(q g \rightarrow q g) & =\frac{4 \pi \alpha_s^2}{9 \hat{s}^2}\left[-\frac{\hat{u}}{\hat{s}}-\frac{\hat{s}}{\hat{u}}+\frac{9}{4}\left(\frac{\hat{u}^2+\hat{s}^2}{\hat{t}^2}\right)\right], \\
		\frac{d \sigma}{d \hat{t}}(g g \rightarrow q \bar{q}) & =\frac{\pi \alpha_s^2[\hat{\hat{u}}}{6 \hat{s}^2}\left[\frac{\hat{t}}{\hat{t}}-\frac{9}{\hat{u}}\left(\frac{\hat{u}^2+\hat{t}^2}{4}\right)\right], \\
		\frac{d \sigma}{d \hat{s}}(g g \rightarrow g g) & =\frac{9 \pi \alpha_s^2}{2 \hat{s}^2}\left[3-\frac{\hat{t} \hat{u}}{\hat{s}^2}-\frac{\hat{s} \hat{u}}{\hat{t}^2}-\frac{\hat{s} \hat{t}}{\hat{u}^2}\right] .
	\end{aligned}
\end{equation}
然后，基于以上所有的数图贡献，我们可以得到领头阶下单举强子产生的散射截面:
\begin{equation}
	\begin{array}{r}
		E \frac{d^3 \sigma}{d^3 p}=\sum_{a b c d} \int_{x_{a, \min }}^1 d x_a \int_{x_{b, \min }}^1 d x_b f_{a / A}\left(x_a\right) f_{b / B}\left(x_b\right) D_{h / c}\left(z_c\right) \\
		\times \frac{1}{\pi z_c} \frac{d \sigma}{d t}(a b \rightarrow c d)
	\end{array}
	\label{eq:CS_LO}
\end{equation}
其中，各个动量分数的积分下限分别为：$x_{a,min}=(x_T e^y)/(2- x_T e^y)$、$x_{b,min}=(x_a x_T e^{-y})/(2x_a - x_T e^y)$ 和 $z_c = (x_T/2x_b)e^{0y} + (x_T/2x_a)e^y$，其中有快度$y$和末态横动量分数$x_T$。
对于次领头阶的计算，需要同时考虑 $2\to 2$ 过程的辐射修正和 $2\to 3$ 过程~\ref{Harris:2001sx,Kidonakis:2000gi,Jager:2002xm}。

	在pQCD次领头阶对部分子模型的修正中，单举强子产生的不变散射截面的表达式与公式~\ref{fig:CS_LO}相似，即为核中部分子分布函数（nPDF）、次领头阶下的部分子散射截面$d\sigma / d \hat{t} (ab \to cd)$和部分子碎裂函数$D_{h/c}(z_c,Q^2)$三部分的卷积。在经过一系列的计算后简化为~\parencite{Kidonakis:2000gi}：
\begin{equation}
	\begin{array}{r}
		\frac{1}{p_T} \frac{d \sigma}{d p_T}= \int F_q (\frac{p_T}{z_h}) D_{q \to h} (z_h, Q^2) \frac{dz_h}{z_h^2} \\
		+ \int F_g (\frac{p_T}{z_h}) D_{g \to h} (z_h, Q^2) \frac{dz_h}{z_h^2}
	\end{array}
	\label{eq:CS_parton}
\end{equation}
其中，$z_h$是碎裂标度$Q$下末态强子携带的动量占母部分子动量的分数；$F_{q,g}(F_T/z_h)$是PDF和$d\sigma / d \hat{t} (ab \to cd)$的卷积。对于PDF的不同参数货形式的如CT18NLO~\parencite{Hou:2019efy}和CT14NLO~\parencite{Hou:2016sho,Dulat:2015mca}。这些不同的nPDF对高横动量末态强子的产生影响都很小。


% Catani and Seymour
%
具体来说，对次领头阶计算的主要困难是奇异性。一般来说，在计算高阶QCD散射截面时，需要考虑实发散贡献和虚修正，并且需要处理不同类型的奇异性。紫外奇异性存在于虚修正中，并通过重整化来消除。红外奇异性存在于虚修正和实发散贡献中，但它们的总和是有限的。
有几种通用的方法可用于抵消红外发散，可用于建立蒙特卡罗程序~\parencite{Nagy:1996bz,}。
在计算次领头阶修正时，以Catani-Seymour偶极子减除方法为例~\parencite{Catani:1996vz}作为介绍。

对于具有两个初态强子的过程，次领头阶散射截面是入射强子的部分子密度函数和硬散射横截面的卷积：
\begin{equation}
	\begin{aligned}
		\sigma(p, \bar{p}) \equiv & \sum_{a, b} \int_0^1 d \eta f_{a / H}\left(\eta, \mu_F^2\right) \int_0^1 d \bar{\eta} f_{b / \bar{H}}\left(\bar{\eta}, \mu_F^2\right) \\
		& \times\left[\sigma_{a b}^{L O}(\eta p, \bar{\eta} \bar{p})+\sigma_{a b}^{N L O}(\eta p, \bar{\eta} \bar{p})\right]
			\label{eq:cs1}
	\end{aligned}
\end{equation}
其中$f_{i/H}$是因子化标度为$\mu_F$时，入射强子$H$里味道为$i$的部分子密度。相应的部分子层次的散射截面为：
\begin{equation}
	\sigma_{a b}^{L O}(p, \bar{p}) \equiv \int_3 d \sigma_{a b}^B(p, \bar{p})=\int_3 d \Gamma^{(3)}\left\langle\left|\mathcal{M}_{a b}^{(3)}\right|^2\right\rangle F_J^{(3)}
		\label{eq:cs2}
\end{equation}
次领头阶的修正为：
\begin{equation}
	\begin{aligned}
		\sigma_{a b}^{N L O}(p, \bar{p}) \equiv & \int d \sigma_{a b}^{N L O}(p, \bar{p})=\int_4 d \sigma_{a b}^R(p, \bar{p}) \\
		& +\int_3 d \sigma_{a b}^V(p, \bar{p})+\int_3 d \sigma_{a b}^C(p, \bar{p})
	\end{aligned}
			\label{eq:cs3}
\end{equation}
其中$d\sigma^R$ 和$d\sigma^V$分别是对散射截面的实贡献和虚贡献。对$d\sigma^C$表示共线减除相。
这里的有限项依赖因子化标度。$d\sigma^C$表达式有限项部分的实际形式定义了因子化标度。

在我们的计算中，我们选择了常用的修正最小减除方案。部分子密度函数也依赖于此方案，因此在公式~\ref{eq:cs1}中的强子横截面中这种依赖关系被抵消掉了。
公式~\ref{eq:cs3}右侧的三个积分分别发散，但它们的和是有限的，由喷注函数$F_J^{m}$提供。喷注函数定义了一个红外安全的量，形式上意味着
\begin{equation}
	F_J^{(4)} \to F_J^{(3)}
\end{equation}
这是在任何四夸克和三夸克构型在动力学上具有相同的运动学特征的情况下。
奇点的存在意味着需要对各个部分进行正则化并取消发散。










%------------------------------- 2.2
\section{喷注的重建算法}
	高能核物理中的喷注重建算法在研究重离子碰撞实验中扮演着重要的角色。喷注是由高能碰撞中产生的高能量和高动量的粒子束，通常由夸克和胶子组成。实验上的探测器只能收集到强子化后的末态粒子信息。但是我们可以通过喷注重建算法将一定锥角内的粒子加以限制来重建一个整体喷注，如图~\ref{fig:jet_reconst}。

  \hspace*{\fill}
\begin{figure}[htbp]
	\centering
	\includegraphics[width=2.3in]{jet_reconst}
	\caption{重建的整体喷注示意图~\parencite{Salam:2007xv}。}
	\label{fig:jet_reconst}
\end{figure}

	目前主要的喷注重建算法分为锥角算法（Iterative Cone Algorithm）~\parencite{Blazey:2000qt}和连续组合算法~\parencite{Ellis:1993tq}两大类。
	其中前者包括Snowmass算法和SIS锥角算法~\parencite{Salam:2007xv}；
	后者包括Cambrige-Aachen算法~\parencite{Dokshitzer:1997in,Wobisch:1998wt}、$k_T$算法~\parencite{Catani:1993hr}和anti$-k_T$算法~\parencite{Cacciari:2008gp}。
	如图~\ref{fig:anti_kt}展示了其中四种不同算法重建的部分子层面喷注示意图：SIS锥角算法、Cambrige-Aachen算法、$k_T$算法和anti$-k_T$算法。
	这些算法在处理背景扣除、强子化等非微扰过程时并不相同，但都保证了红外安全和共线安全（IRC）。
	喷注重建算法的设计和优化涉及到多个因素，包括探测器分辨率、背景噪声、喷注形状和喷注内部结构等。研究人员通过使用模拟数据和实验数据进行验证和优化，以确保喷注重建算法在实际应用中具有高效性和准确性。
	喷注重建算法在高能核物理实验中具有广泛的应用，例如胶子物理~\parencite{JADE:1979rke,PLUTO:1979dxn,Barber:1979yr,TASSO:1979zyf}、顶夸克物理质~\parencite{D0:1995jca,CDF:1995wbb}以及寻找新的粒子和物理现象质~\parencite{CMS:2018uxb,ATLAS:2018mme,ATLAS:2018kot}等。
	部分子穿越QGP前后会发生变化，通过精确的重建喷注，以喷注观测量作为探针来研究QGP的方案成为喷注层析~\parencite{Wang:2002ri}。
	研究人员以此获得关于高能碰撞中发生的物理过程和粒子产生机制的重要信息，从而推动对核物理的深入理解和探索。
	
  \hspace*{\fill}
\begin{figure}[htbp]
	\centering
	\includegraphics[width=4.6in]{anti_kt}
	\caption{由四种不同的喷注算法重建的整体喷注效果示意图~\parencite{Cacciari:2008gp}。}
	\label{fig:anti_kt}
\end{figure}

	几种算法各有优劣和侧重点。接下来，进行简要介绍。
	\textbf{SIS锥角算法}
		这一算法的思想是在$\eta - \phi$的粒子平面上，在以喷注种子为中心的锥角$R$内来重建喷注。首先对横动量高于阈值的粒子依次排序来作为喷注的种子。其次在种子中选取横动量最大的种子$i$，计算其与其它粒子的距离$d_{ij} = \sqrt{(\Delta \phi)^2 + (\Delta \eta)^2}$，如果所有$d_{ij} > R$，那么该种子即重建为喷注；如果$d_{ij} > R$，则该粒子与种子合并为新种子（四动量相加），然后对新种子重复以上步奏，直到无合适的粒子，即得到了稳定的喷注，此时对已经重建为喷注的粒子从列表移除。最后对剩余种子继续重复第二部操作。最后所有种子将重建为喷注。这种算法较简单且效率高，但存在红外安全的问题。

	\textbf{连续组合算法}
		任意两个入射粒子$i$和$j$间距离为$d_{ij}$，入射粒子$i$与粒子束$B$间的距离为$d_{iB}$，分别定义为：
\begin{equation}
\begin{aligned}
	& d_{i j}=\min \left(k_{T i}^{2 p}, k_{T j}^{2 p}\right) \frac{\Delta R_{i j}^2}{R^2}, \\
	& d_{i B}=k_{T i}^{2 p}
\end{aligned}
\end{equation}

	其中喷注半径$R$有 $\Delta R_{ij}^2 = \left(y_i-y_j\right)^2+\left(\phi_i-\phi_j\right)^2$；$k_{Ti}$ 、$y_i$和$\phi_i$分别为粒子$i$的横动量、快度和方位角；参数$p$的不同取值对应了不同算法：$p=0$、$p=1$和$p=-1$分别为Cambrige-Aachen算法、$k_T$算法和anti$-k_T$算法。
	
	Cambrige-Aachen算法：
\begin{equation}
	\begin{aligned}
		& d_{i j}=\frac{R_{i j}^2}{R^2} \\
		& d_{i B}=1
	\end{aligned}
\end{equation}

	$k_T$算法：
\begin{equation}
	\begin{aligned}
		d_{i j} & =\min \left(k_{T i}^2, k_{T j}^2\right) \frac{\Delta R_{i j}^2}{R^2} \\
		d_{i B} & =k_{T i}^2 .
	\end{aligned}
\end{equation}

	anti$-k_T$算法：
\begin{equation}
	\begin{aligned}
		& d_{i j}=\min \left(k_{T i}^2, k_{T j}^2\right) \frac{\Delta R_{i j}^{-2}}{R^{-2}} \\
		& d_{i B}=k_{T i}^{-2}
	\end{aligned}
\end{equation}

此算法由满足条件$d_{iB} < d_{ij}$ 的最小$d_{iB}$开始，将粒子$i$确定为末态喷注并移除重建序列；若满足条件件$d_{iB} > d_{ij}$ ,则将粒子$i$和$j$的四动量相加构建成喷注的一部分，重复这个过程直到符合条件的喷组组分再无剩余粒子。
此方法能保证红外安全和贡献安全，但硬粒子会影响喷注的形状，所以这样重建的结果会出现不规则形状的喷注。此外每次比较的方法也加大了计算时长。
对于Cambrige-Aachen算法，因此时粒子横动量的指数为0，故不受其影响；对于$k_T$算法，重建后的喷注包含较多横动量偏低的粒子，导致喷注几何形状很不规则；而anti$-k_T$算法使粒子间距反比于横动量的平方，重建过程就围绕高横动量的粒子进行，即受软辐射的影响较小，得到的喷注形状较为规则。

% From old chapter 3
%------------------------------- 2.3（2.2.1） ？？？
\section{质子-质子碰撞中的双喷注的产生}
\label{sec:pp}

% PRD II A
	高能核物理中的一个重要的过程就是质子-质子（$pp$）碰撞中双喷注的产生。具体而言，在部分子层次上，这个过程与由两个相撞质子中的部分子 $a$ 和 $b$ 引发 $N$ 个部分子的产生过程有关。 其中 $a$ 和 $b$ 可以是夸克或胶子且$N\geq2$。
	在微扰量子色动力学中，$pp$ 碰撞产生 $N$ 个粒子的散射截面通常可以表示为部分子分布函数 $f_i(x,Q^2)$ 和硬的 $2\!\rightarrow\!N$ 过程的散射截面  $d\hat{\sigma}^{[2\rightarrow N]}_{ab}$的卷积~\parencite{Gao:2012he} ：
\begin{align}
  \frac{d\sigma}{d\Phi_N}=\sum_{a,b}\int_0^1 \!\!\!dx_a\!\int_0^1 \!\!\!dx_b
  \,f_a(x_a,\mu^2_f)\,f_b(x_b,\mu^2_f)\cr
  \times\frac{\,d\hat{\sigma}^{[2\rightarrow N]}_{ab}(x_a,x_b,\mu_r,\mu_f)\,}{d\Phi_N},
  \label{eq:partonic}
\end{align}
其中 $\Phi_N\equiv\Phi_N(p_1,\ldots,p_N)$ 表示末态 $N$ 个粒子的相空间，$x_a$($x_b$) 是来自前向(后向)质子的入射粒子所携带的动量分数，而 $\mu_f$($\mu_r$) 则是因子化(重整化)标度。方程(\ref{eq:partonic})的右侧对粒子 $a$ 和 $b$ 的味道求和，基本上包括夸克-夸克、夸克-胶子和胶子-胶子的初始态。

在微扰计算的次领头阶中，只涉及到 $2\rightarrow2$ 和 $2\rightarrow3$ 的部分子过程，这时双喷注的观测量，比如由双喷注的运动学变量 $V^{(m)}=\{v_1,\ldots,v_m\}$ 所定义的 $m$ 重微分散射截面 ${d\sigma}/{\prod_{j=1}^{m} dv_j}$，可以使用下面公式进行计算~\parencite{Gao:2012he}
\begin{align}
  \frac{d\sigma}{\prod_{j=1}^{m} dv_j}\!=\!\int d\Phi_2(p_1,p_2)\,\frac{d\sigma}{d\Phi_2}\,S_2(p_1,p_2)
  \,\,\,\,\,\,\,\,\,\,\,\,\,\,\,\,\,\,\,\,\,\cr
  +\!\int d\Phi_3(p_1,p_2,p_3)\,\frac{d\sigma}{d\Phi_3}\,S_2(p_1,p_2,p_3).
  \label{eq:NLO}
\end{align}
这里对末态所施加的所有约束条件，包括喷注算法，都融入到函数 $S_N(p_1,\ldots,p_N)$ 中，这些函数应该被适当地定义以满足红外安全条件~\parencite{Gao:2012he}。

在因子化的框架下，方程 (\ref{eq:partonic}-\ref{eq:NLO}) 将可测量的双喷注的散射截面与不可直接测量的部分子分布函数 $f_i(x,Q^2)$ 联系起来。通常，为了在可观测的层面上揭示部分子分布函数对 $x$ 和 $Q^2$ 的依赖关系，可以通过适当选择双喷注变量 $V^{(m)}$，在领头阶（LO）近似下建立 $V^{(m)}$ 的值 $x$ 和 $Q^2$ 值之间的一一对应关系。

例如，可用于定义双重微分散射截面的一个适用的变量组是 $V^{(2)}=\{M_{{\!J\!J}},y_{\textrm{dijet}}\}$~\parencite{Eskola:2013aya}，其中包括喷注对的不变质量的平方 $M^2_{{\!J\!J}}\equiv(p_1+p_2)^2$ 和快度 $y_{\textrm{dijet}}$。需要注意的是，在DY测量中通常使用类似的选择。通过这种选择，我们就可以通过领头阶的关系将 $V^{(m)}$ 与 $x_{a(b)}$ 联系起来：
\begin{align}
  x_{a(b)}=\frac{M_{{\!J\!J}}}{\sqrt{s}}e^{\pm y_{\textrm{dijet}}},
  \label{eq:my}
\end{align}
这样就可以实现对 $x_{a(b)}$ 的动力学扫描。除此之外，还有一个物理标度 $M_{{\!J\!J}}$ 得到了很好的控制。
与DY过程不同，DY过程的典型硬尺度是轻子对的不变质量，而在双喷注产生中有一个额外的物理尺度，即喷注横向动量(~($p_{T1}\!=\!p_{T2}$ 在 LO时))，这在上述 $V^{(m)}$ 选择中尚未完全得到控制。

一般而言，一旦考虑到高阶修正，即考虑到更多的物理效应，已经建立的领头阶对应关系可能会被或多或少地打破。因此，对于实际的可观测量，与各种 $x$, $Q^2$ 和 $i$（味道）值相关的子过程会共同对单个 $V^{(m)}$ 的测量做出贡献，这不可避免。下面从运动学角度给出一个示例。

\begin{figure}[htbp]
	\centering
	\includegraphics[width=2.7in]{2-n}
	\caption{质子-原子核碰撞的硬散射过程的一般性示意图。展示了初始两个部分子碰撞，末态产生n个粒子。}
	\label{fig:2-n}
\end{figure}

对于N个部分子的产生过程,如图~\ref{fig:2-n}，可以通过将所有 $N$ 个末态部分子的光锥动量的前向（后向）分量相加，来推断初始状态的动量分数 $x_a$~($x_b$)~\parencite{Ellis:1994dg}：
\begin{align}
  x_{a(b)}=\sum_{n=1}^N \frac{p_n^{\pm}}{\sqrt{s}}=\sum_{n=1}^N \frac{E_{Tn}}{\sqrt{s}}e^{\pm y_n},
  \label{eq:inferx}
\end{align}
其中，光锥分量为 $p_n^{\pm}\equiv p^0_n\pm p^3_n=E_{Tn}e^{\pm y_n}$，$E_{Tn}=(m^2+p_{Tn}^2)^{1/2}$ 为横向能量，第 $n$个部分子的快度为$y_n$。公式说明了把末态所有粒子的光锥动量的向前分量求和除以根$s$即为$x_a$；把末态所有粒子的光锥动量的向后分量求和除以根$s$即为$x_b$。由此可以看到，当掌握了所有末态信息时，我们就可以准确的反推出初态的动量分数。
然而，在双喷注的非遍举测量中，由于位于双喷注锥之外的末态粒子不受限制，因此不能完全推断出 $x_a$ 和 $x_b$ 的值，这种情况可能发生在NLO及其以上的阶次中。这些问题也说明了全局QCD分析对于部分分布函数的必要性和复杂性。

尽管如此，由于LO过程会给出主导贡献，因此我们仍然可以期望双喷注的三重微分散射截面对于部分分布函数的 $x$ 和 $Q^2$ 依赖依然具有很好的分辨能力，这对于全局分析是有益的~\parencite{CMS:2017jfq}。

	在本章的工作中，我们计算了在微扰QCD框架里次领头阶贡献下双喷注的双重微分散射截面。
	我们使用的程序是 NLOJet++。
	用于预测发现单个喷注或喷注对的概率，其中喷注的产生伴随着任意末态 $p + p \to \text{jet} + X$ 和 $p+p \to \text{jet} + \text{jet} + X$。
 	NLOJet++程序可被用于计算散射截面的不同阶次（领头阶和次领头阶），其核心库使用了Catani-Seymour偶极子减除方法，详见\ref{sec:nlojet_cs}节所介绍，并提供了一个通用框架供各种过程的计算模块使用，从而促进了对高能碰撞过程的理论研究和实验数据的比较分析~\parencite{Nagy:2001fj,Nagy:2003tz}。

为了测试数值计算，我们首先在图~\ref{fig:eta_cms_pp} 中展示了在质心系能量为13 TeV的质子-质子碰撞中测量得到的双喷注赝快度~[$\eta_{\textrm{dijet}}\!=\!(\eta_1\!+\!\eta_2)/2$] 分布，该数据由CMS合作组测量得到~\parencite{CMS:2018jpl}。
图中 $\eta_{\textrm{dijet}}$ 是在 $p$Pb 碰撞的实验室参考系中定义的，因此与质心系中的 $\eta_{\textrm{dijet}}=\eta^{\textrm{cm}}_{\textrm{dijet}}+0.465$~\parencite{CMS:2018jpl}有关系。理论结果分别在CT14 领头阶和次领头阶质子部分子分布函数下进行了领领头阶（虚线）和次领头阶（实线）计算。图中的圆圈代表CMS数据（垂直误差条很小）~\parencite{CMS:2018jpl}。
喷注采用 anti-$k_T$ 算法，并使用锥尺寸 $R=0.3$ 以及实验室参考系下的快度截断 $|y_{\textrm{lab}}|<3$。
双喷注的平均横向动量被限制在 $75<p_{T,\textrm{avg}}<95$ GeV 的范围内。
喷注对的相对方位角被限制在$|\Delta\phi_{12}|>2\pi/3$。
在该分析中，对于领头喷注和次领头喷注的横向动量施加了截断，分别为$p_{T1}>30$ GeV和$p_{T2}>20$ GeV~\parencite{CMS:2018jpl}。
我们发现，使用CT14质子的部分子分布函数~\parencite{Dulat:2015mca} 和取因子化和重整化尺度为 $\mu_0=\mu_f=\mu_r=p_{T,\textrm{avg}}$ 的NLO计算，可以在一定程度上较好地描述CMS实验数据（详见文献~\parencite{Eskola:2019dui}的$\chi^2$分析）。

\hspace*{\fill}
\begin{figure}[htbp]
  \centering
  \includegraphics[width=3.3in]{eta_cms_pp}
  \caption{质心系能量为 $\sqrt{s}=5.02$~TeV 的质子-质子碰撞中，双喷注赝快度分布的归一化结果。}
  \label{fig:eta_cms_pp}
\end{figure}
 
由于 $\eta_{\textrm{dijet}}$ 分布是在双喷注的平均横向动量 $p_{T,\textrm{avg}}$ 被限制的情况下测量得到的，因此图~\ref{fig:eta_cms_pp} 中的归一化分布 $N^{-1}dN/d\eta_{\textrm{dijet}}$ 实际上对应于带有$V^{(2)}=\{\eta_{\textrm{dijet}},p_{T,\textrm{avg}}\}$ 的双重微分截面 $\sigma^{-1}d^2$ $\sigma/d\eta_{\textrm{dijet}}dp_{T,\textrm{avg}}$。
需要注意的是，即使在LO阶段，使用特定的 $\{\eta_{\textrm{dijet}},p_{T,\textrm{avg}}\}$ 集合仍无法唯一确定双喷注不变质量 $M_{{\!J\!J}}$ 或 $x_{a(b)}$，而比值 $x_a/x_b=e^{2\eta_{\textrm{dijet}}}$ 是固定的。总体而言，随着 $\eta_{\textrm{dijet}}$的增加， $x_{a(b)}$ 的平均值也会相应增加（减少）~\parencite{CMS:2018jpl}。我们将在滴~\ref{sec:triple}章中进一步讨论这个问题。



%------------------------------- 2.3
\section{高能碰撞中的蒙特卡洛事件产生器}

	在粒子物理学中，两个人射粒子碰撞的产物或者粒子的独立衰变过程被称作一个“事件（event）”。
	出射粒子的数量和它们的性质由于量子过程的随机性，在每个事件中都会有所变化。
	如果给出一个能够连续生成这种模拟事件随机序列的数值算法，就有可能产生与数据相比较的模拟事件集合，此即“事件生成器（event generator）”，它是利用计算机在基本粒子之间的高能碰撞中生成尽可能详细事件的标准工具，就像一个完美探测器所观测到的那样。
	这些模拟可以基于已知的或假设的自然法则。这样就可以探索和比较不同的模式，并研究所提出的物理可观测量对这些差异的敏感性。
	简单来说，高能物理事件的结构是碰撞中产生的亚原子粒子列表（称为“事件记录（event record）”，包括每个粒子的四动量、产生点、显示事件建模的中间阶段的历史信息和许多其他属性），以及该事件发生的概率的度量（由该事件的权重相对于样本的权重总和所给出）。
	事件生成器可以预测高能碰撞或衰变过程的所有可观测性质。然而，目前无法仅凭第一原理计算出事件的全部属性。许多不同的复杂现象可能彼此相关，并由众多模型描述，每个模型都专注于有限的动力学范围。
	目前领域内常用的蒙特卡洛（Monte-Carlo，简称MC） 事件产生器有 PYTHIA~\parencite{Sjostrand:2006za}, SHERPA~\parencite{Gleisberg:2008ta}，POWHEG~\parencite{Frixione:2007vw,Alioli:2010xd}, HERWIG~\parencite{Bahr:2008pv}等，其中SHERPA和POWHEG包含了对次领头阶散射截面的高阶修正。其中的一部分物理是严格从理论中推导出来的，而另一部分则基于现象学模型，其中的参数需要从数据中确定。
	通常都能给出事件中的完整信息，比如硬过程（hard process）、潜在事件(underlying event)、部分子簇射(parton shower）、多重散射（multiparton interactions），碰撞残余粒子（beam remnants）和强子化（hadronization）过程等。
	如图~\ref{fig:pythia}，给出了用PYTHIA8 完全模拟pp 碰撞中$\to t\bar{t}$事件的结构示意。

  \hspace*{\fill}
	\begin{figure}[htbp]
	\centering
	\includegraphics[width=4.2in]{pythia}
	\caption{PYTHIA8 完全模拟的pp$\to t\bar{t}$事件结构示意图。径向坐标表示硬度标度，从中心附近最硬的子过程（标记为 ~d$\hat{\sigma_0}$）开始，到最外围稳定的末态粒子和入射粒子束~\parencite{Bierlich:2022pfr}。}
	\label{fig:pythia}
	\end{figure}

%------------------------------- 2.4

\section{Glauber模型：重离子碰撞的初始几何}

  Glauber模型是为了解决复合粒子的高能量散射问题而开发的，是一种常用于描述重离子碰撞初始几何结构的模型。模型由Glauber于20世纪50年代提出~\parencite{Miller:2007ri}，该理论最初用于描述光学中的散射现象，后来被推广应用于核物理中的碰撞过程。
	在Glauber模型中，重离子被视为由核子组成的多体系统。这些核子之间通过核力或强相互作用相互作用，形成了初始的核物质分布。Glauber模型假设碰撞参与的核子是独立的~\parencite{wong1994introduction}，并且它们在碰撞过程中保持不变。它基于核物质的几何形状和碰撞参数，提供了一种估量碰撞过程中核子相互作用的方法。
	模型的核心思想是通过计算在碰撞中两个核子的相互作用，来确定碰撞过程中的核子核密度分布。具体而言，Glauber模型使用核子-核子碰撞的散射振幅来计算核子的碰撞概率。通过考虑碰撞参与核子的位置和形状，可以得到碰撞过程中的核子数密度分布。
	Glauber模型对于重离子碰撞的初始几何结构提供了一些重要的信息。例如，它可以用于计算碰撞中的核子碰撞数目、核子密度分布的均匀性以及碰撞中的碰撞参数$b$（impact parameter）。这些参数对于理解碰撞过程的能量传递、碰撞产物的分布和相关的物理现象非常关键。
	需要注意的是，Glauber模型是一个简化的模型，它并未考虑诸如QCD效应和核子的内部结构等复杂的物理过程。然而，它在研究重离子碰撞的初步几何结构时仍然具有重要的应用价值，并且为进一步的分析提供了有用的起点。
	
	在使用Glauber方法进行几何参数的计算时，需要提供一些实验数据作为模型输入。其中最重要的是通过低能电子散射实验测得的核电荷密度$\rho(r)$（Nuclear charge densities）以及非弹性核子-核子散射截面$\sigma_{inel}^{NN}$的能量依赖性。
	其中核子密度常通过具有三个参数的费米分布进行参数化：
 \begin{equation}
 	\rho(r)=\rho_0 \frac{1+w(r / R)^2}{1+\exp \left(\frac{r-R}{a}\right)},
 \end{equation}
	这里，$R$为核半径，$a$和$w$分别为核与球形的偏离程度，核中心的核子密度为$\rho_0$。
	有归一化形式：
 \begin{equation}
	\int d^3 r \rho_A (r)=A.	
\end{equation}
	部分原子核的三个参数取值见表~\ref{fig:Glauber_table}。

	\begin{figure}[htbp]
		\centering
		\includegraphics[width=2.4in]{Glauber_table}
		\caption{Glauber模型中氙核、金核和铅核的核子密度费米分布的三参数值~\parencite{Zakharov:2018ctz,Moller:2015fba}。}
		\label{fig:Glauber_table}
	\end{figure}

	两个重离子，目标核A和入射核B，以相对论速度相互碰撞，见图~\ref{fig:Glauber1}。碰撞参数为$\vec{b}$，代表了两个圆盘的对心程度，假设核中核子均匀分布的话，碰撞参数写成标量$b$。若$b\gg 1$，则表示边缘（peripheral）碰撞，若$b\sim 0 $，则表示中心（central）碰撞，如图~\ref{fig:Glauber2}。
	我们关注位于与目标核心的中心相对位移$\vec{s}$以及偏离入射核心的中心距离$\vec{s} - \vec{b}$。
	
  \hspace*{\fill}
\begin{figure}[htbp]
	\centering
	\includegraphics[width=3.5in]{Glauber1}
	\caption{Glauber模型中核-核碰撞示意图，其中碰撞参数为 $b$。左图（侧视图，即垂直于束流方向）为两个重离子在碰撞之前，；右图（平行于束流方向）~\parencite{Miller:2007ri}。}
	\label{fig:Glauber1}
\end{figure}

\begin{figure}[htbp]
	\centering
	\includegraphics[width=4.2in]{Glauber2}
	\caption{横向平面上的三种重离子碰撞（从左到右）：边缘碰撞、半中心碰撞和中心碰撞。}
	\label{fig:Glauber2}
\end{figure}

	定义厚度函数$T_A(s)$为一个给定核A中的核子沿$z$轴方向的投影：
\begin{equation}
	T_A(s)=\int \rho_A(r) d z, \quad \int d^2 s T_A(s)=A
\end{equation}
	重叠函数为$T_A(s)$和$T_B(s)$的卷积：
\begin{equation}
	T_{A B}(b)=\int d^2 s T_A\left(s-\frac{b}{2}\right) T_B\left(s+\frac{b}{2}\right), \int d^2 s T_{A B}(b)=A B
\end{equation}
碰撞发生的概率为$T(b) \sigma_{inel}^{NN}$，弹性过程几乎没有能量损失，因此在Glauber模型的计算中通常不考虑。于是，发生一次二体碰撞的概率为：
\begin{equation}
	P_{A B}^{(1)}(b)=\frac{T_{A B}(b)}{A B} \sigma_{i n e l}^{\mathrm{NN}}, \quad \int d^2 b P_{A B}^{(1)}(b)=\sigma_{i n e l}^{\mathrm{NN}}
\end{equation}
	核A、B至少发生一次二体碰撞的概率遵循柏松分布：
\begin{equation}
	P_{\text {inel }}^{A B}(b) \approx 1-\exp \left(-T_{A B}(b) \sigma_{\text {inel }}^{\mathrm{NN}}\right)
\end{equation}
	于是，总散射截面即碰撞概率对碰撞参数$b$对积分：
\begin{equation}
	\sigma_{\text {inel }}^{A B}=\int d^2 b P_{\text {inel }}^{A B}(b)=\int d^2 b\left(1-\exp \left(-T_{A B}(b) \sigma_{\text {inel }}^{\mathrm{NN}}\right)\right)
\end{equation}
另外两个重要参数即核子-核子二体碰撞总数和参与非弹性碰撞的核子总数，分别为：
\begin{equation}
		N_{\text {coll }}^{A B}(b) =A B N_{\text {coll }}(b)
\end{equation}
\begin{equation}
	\begin{aligned}
		N_{\text {part }}(b) & =A \int T_A(s+b / 2) \frac{d \sigma_{\text {inel }}^B}{d s^2}+B \int T_B(s-b / 2) \frac{d \sigma_{\text {inel }}^A}{d s^2} d s^2 \\
		& =A \int T_A(s+b / 2)\left\{1-\left[1-T_B \sigma_{\text {inel }}^{N N}\right]^B\right\} d s^2 \\
		& +B \int T_B(s+b / 2)\left\{1-\left[1-T_A \sigma_{\text {inel }}^{N N}\right]^A\right\} d s^2
	\end{aligned}
\end{equation}


%------------------------------- 2.5
\section{部分子分布函数及其核修正}
%\section{冷核物质效应}

	人们对于强相互作用在高能行为方面的现代理解源于上世纪60年代末的部分子模型构想。随着上世纪70年代初渐进自由的出现，这一理解得到确认。部分子模型的提出旨在解释SLAC所进行的深度非弹性电子散射实验中所观察到的规律~\parencite{Bloom:1969kc,Breidenbach:1969kd}。
	在~\ref{sec:SM}节中已提到每个核子由三个上或下价夸克所组成，外加虚胶子、海夸克和反夸克。其中核子内的胶子、海夸克和反夸克是短暂存在，时隐时现的~\parencite{ellis2003qcd}。夸克和胶子的定义基于部分子（夸克或胶子）分布函数$q(x,Q^2)$和$g(x,Q^2)$，这提供了部分子在核子中的概率分布，其中$x$是Bjorken动量分数，即该部分子所携带的核子动量的比例，它是在无限动量参考系中测量的，$Q^2$是进行测量的动量标度。

	涉及原子核的高能碰撞中的硬过程无论是在实验条件下还是在自然条件下发生，准确描述这些过程都需要对部分子分布函数有所了解。
	原子核部分子分布函数是描述原子核内部部分子（如夸克和胶子）的分布情况的函数。

	由于将核子结合在一起的额外动力学，大型原子核的部分子结构显然会偏离自由核子的结构~\parencite{Arneodo:1992wf, Armesto:2006ph, Cloet:2012td, CLAS:2019vsb, Kulagin:2004ie, Ru:2016wfx, Wang:2020uhj}。
	在 $p$A 碰撞中，原子核的部分子分布函数和自由核子的部分子分布函数之间的差异可能导致相对于 $pp$ 碰撞而言的双喷注散射截面的核修正，这被称为初态冷核物质效应（cold nuclear matter effects，简称CNM）~\parencite{Albacete:2017qng}。
	准确确定这种偏差不仅有助于理解束缚核子之间的相互作用，而且还为准确研究其他物理学提供了基准，例如研究相对论核-核碰撞中热密的夸克-胶子等离子体的末态喷注淬火现象~\parencite{Cao:2020wlm, Gyulassy:2003mc}。
	原子核部分子分布与自由质子的部分子分布的偏差通常用量化为比值：
\begin{align}
	r^A_i(x,Q^2)=\frac{f^{A,proton}_i(x,Q^2)}{f^{proton}_i(x,Q^2)}.
	\label{eq:ri}
\end{align}
通过对味道$i$求和以及对动量分数$x$求积分，可以得到归一化形式：
\begin{align}
\sum_i \int  dx f_{i}(x,Q^2)=1
\end{align}
一般来说，核修正因子$r^A_i(x,Q^2)$由核环境所附加的非微扰动力学来决定，并且取决于核质量数$A$、部分子味道类型 $i$、动量分数 Bjorken-$x$和分辨率标度 $Q^2$,如图~\ref{fig:nPDF_R}，可以明显观察到这种依赖性。
$Q^2$越大，得到的图像越详细。例如，在更高的$Q^2$下观测时，在某个$Q^2$下的胶子可能会波动成夸克-反夸克对。因此，夸克和胶子的分布会随着$Q^2$的变化而演化。
\hspace*{\fill}
\begin{figure}[htbp]
	\centering
	\includegraphics[width=3.7in]{nPDF_R}
	\caption{在两个不同标度$Q^=1。69$ GeV和$Q^=100$ GeV下，铅核不同类型的核修正因子$R_V, R_S, R_G$~\parencite{Eskola:2009uj,}。}
	\label{fig:nPDF_R}
\end{figure}

\hspace*{\fill}
\begin{figure}[htbp]
	\centering
	\includegraphics[width=2.8in]{nPDF}
	\caption{核PDF与自由核子PDF之比 $r_i^A$ 作为Bjorken $x$的函数~\parencite{Eskola:2009uj}。}
	\label{fig:nPDF}
\end{figure}
从现象学的角度，核效应可以由$r^A_i(x,Q^2)\ne 1$来量化。如图~\ref{fig:nPDF}所示，
可以观察到束缚核子的PDF与自由核子的不同而从而导致部分子分布函数不同类型的核修正。
在nPDFs的整个运动学区间，这些修正发生在不同的Bjorken-$x$值上。
在小Bjorken-$x$值上，例如小于0.05时，与自由核子相比，原子核内核子的PDF受到削弱，这个过程被称为遮蔽（shadowing）效应~\parencite{Glauber:1955qq, Gribov:1968jf}。这是由于高能量下原子核的洛伦兹收缩和核子的重叠所导致的；
但是，由于动量守恒，如果存在核遮蔽效应，那么在一些中等Bjorken-$x$值上就会存在反遮蔽（antishadowing）效应，即核内的PDF得到增强；
随Bjorken-$x$值的增大，会观察到部分子分布函数的减少，出现EMC（European Muon Collaboration ）效应；
在非常大的Bjorken--$x$处，原子核内的核子由于其量子机制的性质而发生的随机运动，核内的核子没有确定的位置，而是占据由海森堡不确定性原理所决定的一系列动量状态，即费米移动（Fermi motion）效应~\parencite{Arneodo:1992wf,Armesto:2006ph}。

% PRD：Introduction：Fisrt
在研究高能核碰撞中的各种硬散射过程中，PDF是不可或缺的非微扰输入量~\parencite{Paukkunen:2018kmm,Kovarik:2019xvh}，包括了轻子-核碰撞、强子-核碰撞和核-核碰撞。它们碰撞的初态都涉及色中性的强子，而非夸克，所以计算都涉及非微扰的PDF。
% PRD：Introduction：Second...
目前存在许多不同的参数化nPDF集合，比如十几年前的EPS09、EKS98、HKN07、nDS、DSSZ和如今的TUJU19、KSASG20、nNNPDF2.0等，这些集合是基于量子色动力学微扰论的因式分解，通过对实验数据进行全局分析来确定的~\parencite{deFlorian:2011fp,Eskola:2009uj,Eskola:2016oht,Eskola:2021mjl,Hirai:2007sx,Khanpour:2020zyu,Kovarik:2015cma,Kusina:2020dki,AbdulKhalek:2020yuc,Wang:2016mzo,Walt:2019slu,Segarra:2020gtj}。
然而，可以观察到这些nPDF集所给出的因子$r^A_i(x,Q^2)$之间存在差异~\ref{fig:nPDFs_uncer}，并且它们的不确定度仍然相当大，见图~\ref{fig:nPDFs_uncer}。
\begin{figure}[htbp]
	\centering
	\includegraphics[width=5.6in]{nPDFs_diff}
	\caption{图中列举了铅核中四种不同的nPDF参数化形式：EPS09、EKS98、HKN07、EPS08和nDS，其中分为三种不同类型夸克：价夸克、海夸克和胶子~\parencite{AbdulKhalek:2020yuc}。}
	\label{fig:nPDFs_diff}
\end{figure}

\begin{figure}[htbp]
	\centering
	\includegraphics[width=5.9in]{nPDFs_uncer}
	\caption{基于次领头阶拟合的EPPS16、nCTEQ、DSSZ12的核修正~\parencite{Paukkunen:2018kmm}。}
	\label{fig:nPDFs_uncer}
\end{figure}
在LHC运行之前，进行全局提取的主要数据来源是从轻子深度非弹性散射（DIS）和强子-核碰撞中的坠尔-颜（DY）过程中获得的数据。
在这两种过程中，通过控制最终状态的运动学变量，可以对$x$和$Q^2$上的nPDF进行详细扫描。
例如，在DIS中，传统的两个变量是Bjorken-$x_B$和虚光子，而在DY中，通常选择双轻子的不变质量$m_{jj}$和快度$y$作为两个变量~\parencite{ellis_stirling_webber_1996}。
然而，DIS和DY过程都对夸克分布非常敏感，但对胶子分布提供的约束较少。因此，对于胶子分布的核修正因子$r^A_i(x,Q^2)$的不确定性通常比对于夸克分布的不确定性更大~\parencite{Eskola:2009uj}。
\begin{figure}[htbp]
	\centering
	\includegraphics[width=3.7in]{nPDF_gluon}
	\caption{利用LHC实验的新数据，针对nCTEQ15和EPPS16这两种不同的nPDF，对铅核中胶子做了核修正~\parencite{Kusina:2017gkz}。}
	\label{fig:nPDF_gluon}
\end{figure}
尽管早期的一些全局分析已经考虑了来自RHIC的$\pi$介子数据，这些数据对核子的胶子分布更为敏感，但碎裂函数可能会在确定nPDF时不可避免的引入额外的理论上的不确定性~\parencite{CMS:2018jpl}。如图~\ref{fig:nPDF_gluon}，是用LHC新产生的实验数据加入胶子的贡献后，对nCTEQ15和EPPS16和做了更精确修正。
来自LHC的丰富实验数据显著促进了对nPDF的提取。例如，通过DY过程产生的弱玻色子，为高分辨率尺度上的原子核的夸克分布提供了新的视角~\parencite{Ru:2016wfx,Paukkunen:2010qg,Ru:2014yma,Ru:2015pfa,Albacete:2013ei}。
尤其是在质子-核（$p$A）碰撞中通过双喷注~\parencite{CMS:2018jpl,Eskola:2019dui}和重味介子~（$J/\Psi$, $D^0$）等的产生~\parencite{Kusina:2017gkz,Kusina:2020dki,Eskola:2019bgf}，可以对原子核的胶子分布提供有力的限制，从而更加巩固了我们对胶子的遮蔽效应和反遮蔽效应的理解~\parencite{Eskola:2019dui,Kusina:2017gkz,Kusina:2020dki,Eskola:2019bgf}。
\begin{figure}[htbp]
	\centering
	\includegraphics[width=4.6in]{nPDF_gluon2}
	\caption{ 在不同的标度$Q^2=10$ GeV$^2$，度$Q^2=10^4$ GeV$^2$下，EPPS16针对铅核中胶子的核修正因子在reweighting前后做了对比~\parencite{Eskola:2019dui}。}
	\label{fig:nPDF_gluon2}
\end{figure}
图~\ref{fig:nPDF_gluon2}中，EPPS16针对铅核中胶子的核修正因子在reweighting后，误差有了很大的降低，并且遮蔽效应和反遮蔽效应有了很大的增强，说明核修正因子对味道确实有重要依赖。
\begin{figure}[htbp]
	\centering
	\includegraphics[width=4.4in]{nPDF_gluon3}
	\caption{ CMS合作组实验关于$p$Pb碰撞相对于$pp$碰撞的比值随双喷注的快度$\eta_{dijet}$的变化结果。这展现在五个不同的$p_T^{ave}$区间内，并将实验结果与两种不同的核修正DSSZ、EPS09以及次领头阶pQCD的计算进行了对比~\parencite{CMS:2018jpl}。}
	\label{fig:nPDF_gluon3}
\end{figure}
在实验上，CMS合作组也考虑了胶子的贡献，如图~\ref{fig:nPDF_gluon3}所示， 在小的$\eta_{dijet}$双喷注区域中主要是EMC效应，在大的$\eta_{dijet}$双喷注区域中主要是遮蔽效应，可以观察到$p$Pb和$pp$数据的比值比单位值1显著偏离，其中DSSZ的EMC效应比EPS09的更弱。


%------------------------------- 2.6
\section{线性玻尔兹曼输运模型}
\label{sec:lbt}	
 
  在高能核碰撞中形成的核物质新状态QGP中，传播的高$p_T$部分子由于与介质相互作用，会同时经历碰撞能量损失和辐射能量损失~\parencite{Baier:2000mf}。
   为了理解和描述高能重离子碰撞中的物质演化过程，我们需要使用动力学模型。
  输运模型为我们研究介质中的物理现象提供了重要的工具和方法。通过使用这些模型，我们可以更好地理解和解释在高能物理实验中观察到的现象，并为理论与实验之间的联系提供重要的支持。

  线性玻尔兹曼输运（Linear Boltzmann Transport，简称LBT）模型是一种当前成熟且广泛使用的描述喷注在热密介质（膨胀的量子色动力学（QCD）介质）中演化的模型~\parencite{Nahrgang:2013saa,Scardina:2017ipo,Cao:2017hhk,Chen:2020tbl,Luo:2018pto,Yang:2022nei,He:2015pra}。该模型包括基于微扰量子色动力学（perturbative QCD）的弹性和非弹性过程。
  其中弹性过程即在热尺度下喷注簇射部分子与介质部分子之间的领头阶散射，
  而非弹性过程即在高扭度（Higher-Twist）能量损失效应形式下介质诱导胶子的辐射过程。
  通过这些散射过程，我们跟踪了相空间演化，涵盖了领头部分子和发射的胶子，以及从QGP背景中散射出的热部分子和介质内留下的能量-动量空洞。
  其中QGP 的动力学演化由 $3 + 1$维 CLVisc流体动力学模型的初始条件给出。
  最初，它是为了研究喷注在介质中以接近光速行进的情况下所产生的所谓马赫锥激发而开发的~\parencite{Li:2010ts}。在这种介质中，声速小于喷注的传播速度~\parencite{Casalderrey-Solana:2004fdk,Stoecker:2004qu,Ruppert:2005uz,Chaudhuri:2005vc}。
  %LBT用于研究喷注在密集QGP介质中的相互作用和传播。
  该模型包含完整的弹性 $2 \to 2$ 散射过程~\parencite{Peng:2014jia}以及具有多个胶子辐射和全局能量动量守恒的非弹性过程 $2 \to 2 + n$~\parencite{Wang:2013cia,Cao:2016gvr,Luo:2016ufs}。
  
  LBT模型的基本构建单元是线性玻尔兹曼方程。
  \begin{equation}
  	\begin{aligned}
  		p_a \cdot \partial f_a= & \int \sum_{b c d} \prod_{i=b, c, d} \frac{d^3 p_i}{2 E_i(2 \pi)^3}\left(f_c f_d-f_a f_b\right)\left|\mathcal{M}_{a b \rightarrow c d}\right|^2 \\
  		& \times \frac{\gamma_b}{2} S_2(\hat{s}, \hat{t}, \hat{u})(2 \pi)^4 \delta^4\left(p_a+p_b-p_c-p_d\right)+\text { inelastic }
  	\end{aligned}
  \end{equation}
	其中的求和遍历所有可能的夸克味道和散射通道。
	$f_i=(2 \pi)^3 \delta^3\left(\vec{p}-\vec{p}_i\right) \delta^3\left(\vec{x}-\vec{x}_i-\vec{v_i}t \right) (i=a, c)$ 是喷注簇射部分子在散射之前和之后以及介质反冲部分子的相空间密度。
	$f_i=1 /\left(e^{p_i \cdot u / T} \pm 1\right)(i=b,d)$是QGP介质中热部分子的相空间分布，其中的局域温度为 $T$。
	流体速度为 $u=(1,\vec{v})/\sqrt{1-\vec{v}}$，
	$\gamma_b$ 是部分子 $b$ 的色-自旋简并度。
  
  领头阶的弹性散射振幅 $\left|\mathcal{M}_{a b \rightarrow c d}\right|^2$ ~\parencite{Eichten:1984eu}具有共线发散，这些发散在LBT模型中通过以下因子进行正规化~\parencite{Auvinen:2009qm}:
  \begin{equation}
  	S_2(\hat{s}, \hat{t}, \hat{u})=\theta\left(\hat{s} \geqslant 2 \mu_D^2\right) \theta\left(-\hat{s}+\mu_D^2 \leqslant \hat{t} \leqslant-\mu_D^2\right)
  \end{equation}
  其中，$\hat{s}, \hat{t}$ 和 $\hat{u}$ 是曼德尔斯通变量。具有3个夸克味道的德拜屏蔽质量为：
  \begin{equation}
  	\mu_D^2=\frac{3}{2} g^2 T^2
  \end{equation}
	相应的弹性散射截面为：
	\begin{equation}
		\frac{\sigma_{ab  \rightarrow cd}}{d \hat{t}} = \frac{\left|\mathcal{M}_{a b \rightarrow c d}\right|^2 }{16\pi \hat{s}^2}
	\end{equation}
		%其中对于末态胶子（夸克），这里忽略了玻色增强以及辐射过程的细致平衡。
		强耦合常数 $\alpha_s = g^2/4\pi$ 是固定的，并将通过拟合实验数据进行调整。
	
	目前的LBT模型仅考虑了弹性散射引起的胶子辐射。胶子辐射的微分非遍举比率遵循高扭度方法中的比率~\parencite{Guo:2000nz,Wang:2001ifa}，则有
  \begin{equation}
  	\frac{d \Gamma_a^{\text {inel }}}{d z d k_{\perp}^2}=\frac{6 \alpha_s P_a(z) k_{\perp}^4}{\pi\left(k_{\perp}^2+z^2 m^2\right)^4} \frac{p \cdot u}{p_0} \hat{q}_a(x) \sin ^2 \frac{\tau-\tau_i}{2 \tau_f}
  \end{equation}
  其中$P_a(z)$ 是传播中的部分子$a$发射一个携带能量分数为$z$的胶子的碎裂函数。横向动量为$k_\perp$，$m$ 是传播部分子的质量。
  $\tau_f = 2p_0z(1-z/(k_{\perp}^2+z^2m^2))$ 是胶子形成时间。$\tau_i$ 是最后的胶子辐射的时间。
  在非弹性过程中，弹性散射比率由喷注输运参数表达：
  \begin{equation}
  	\hat{q}_a(x)=\sum_{b c d} \rho_b(x) \int d \hat{t} q_{\perp}^2 \frac{d \sigma_{a b \rightarrow c d}}{d \hat{t}}
  \end{equation}
  即被定义为平均自由程内的横向动量转移的平方。部分子密度 $\rho_b(x)$ 包括简并因子。上述的碎裂函数$P_z(z)$包含一个红外发散项，并通过德拜屏蔽质量$\mu_D$进行正则化，以作为辐射胶子能量的红外截断。

	在LBT中部分子输运模拟的实际实现中，弹性和非弹性散射的概率在每个小且有限的时间步长$\Delta \tau$ 中一起计算，以确保单位性。在部分子$a$传播过程中的一个时间步长$\Delta \tau$  内进行弹性散射的概率为：
\begin{equation}
		P_{\mathrm{el}}^a=1- \exp \left[-\Delta \tau \Gamma_a^{\mathrm{el}}(x)\right]
\end{equation}
	其中部分子$a$ 的总弹性散射率定义为：
\begin{equation}
		\Gamma_a^{\mathrm{el}} \equiv \frac{p \cdot u}{p_0} \sum_{b c d} \rho_b(x) \sigma_{a b \rightarrow c d}
\end{equation}
	非弹性过程的概率为：
\begin{equation}
	P_{\text {inel }}^a=1-\exp \left[-\Delta \tau \frac{1}{1+\delta_g^a} \int d z d k_{\perp}^2 \frac{d \Gamma_a^{\mathrm{inel}}}{d z d k_{\perp}^2}\right] 
\end{equation}
	其中为 $\Gamma_a^{\mathrm{inel}}(x)$ 为来自部分子$a$的总的胶子辐射率。总散射概率表示为：
\begin{equation}
	P_{tot}^a = P_{el}^a  (1- P_{inel}^a ) + P_{inel}^a 
\end{equation}
	即将其分解为纯弹性散射的概率（第一项）和至少有一个胶子辐射的非弹性散射的概率（第二项）。需要注意的是，对于非常小的时间步长长$\Delta \tau \rightarrow 0$ ，上述单位时间内的总散射概率就是弹性和非弹性散射概率的总和。

	根据弹性散射截面和胶子辐射率的微分形式，对最后部分子、反冲部分子和辐射胶子的散射通道、味道、能量和动量进行采样。在每次散射中，通过多个辐射胶子来确保全局能量和动量的守恒。

	在LBT模型中，每次辐射胶子与介质中的热部分子发生散射时，初始热部分子$b$被记录为“负”粒子，它们也被允许按照玻尔兹曼方程在介质中传播。这些“负”粒子的能量和动量将从所有最终可观测量中减去，以考虑玻尔兹曼输运方程中的反作用。它们是喷注诱导介质激发的一部分，并表现为喷注簇射部分子后方的扩散尾迹~\parencite{Wang:2013cia,Li:2010ts,Peng:2014jia}。
	
%------------------------------- 2.7
\section{SHELL模型}
\label{sec:shell}

  另一个常用的输运模型是 SHELL（Simulating Heavy Quark Energy Loss with Langevin equations）模型。它提供了一种模拟重味部分子在热密介质中演化的方法~\parencite{Svetitsky:1987gq,Cao:2013ita,Dai:2018mhw,Wang:2019xey,Wang:2020ukj,Wang:2020bqz,Wang:2021jgm,Li:2022tcr}。
  在高能核碰撞中，SHELL模型已成功应用于重味喷注的研究，对一系列实验测量结果提供了令人满意的描述。这些实验测量包括重味喷注的$p_T$不平衡~\parencite{Dai:2018mhw}、径向分布~\parencite{Wang:2019xey,Wang:2020ukj,Wang:2020bqz}、碎裂函数~\parencite{Li:2022tcr} 以及$Z^0+$重味强子/喷注的关联~\parencite{Wang:2020qwe,Wang:2021jgm}。
  
  在本节中，我们将详细说明在研究中使用该方程进行模拟重味部分子能量损失的情况。

% 【洒2020 Z0_3_1_15】
  要完整地研究在QGP中传播的重味喷注，需要同时考虑轻味和重味粒子的时空演化~\parencite{Sharma:2009hn,Rapp:2018qla,Djordjevic:2013xoa,Kang:2016ofv}。
  SHELL模型了这些演化并考虑了热密介质内的弹性和非弹性能量损失，接下来我们分别这两种不同的能量损失机制。
  %理论框架。

%张清
 %Langevin 方程是 Foker-Planck 方程的随机实现形式，且不依赖于准粒子（quasi-particle）假设，即 
 %Boltzmann 方程成立的前提是部分子在介质中的能量耗散是与周围谁粒子的散射过程，但 Fokker-Planck~\parencite{} 方程和 Langevin 方程可以直接由介质的谱函数得到输运参数法~\parencite{Liu:2016ysz}。
 

%------------------------------- 2.7.1
\subsection{弹性散射} 
%弹性碰撞能量损失

	重夸克（heavy quark，简称HQ）之所以被认为“重”，首先因为在粒子物理学的背景下，它们的质量比QCD的经典非微扰尺度大得多，即$m_Q \ll \lambda_{QCD}$，这使得我们可以在pQCD中计算它们的产生截面~\parencite{Cacciari:2005rk}。其次，在极端相对论碰撞形成的QCD物质的背景下，它们的质量大于环境介质中典型的温度，即$m_Q \ll T_{QGP}$，这意味着重夸克的产生基本上仅限于初始的硬碰撞（在QGP形成之前的形成时间约为$\sim 1/2m_Q \sim 0.1 $fm/c），此外它们的热化时间与火球寿命相当（甚至更长）~\parencite{Svetitsky:1987gq}。
	在低动量$p_T$下，重夸克的大质量使得我们可以将它们在QGP中的传播视为布朗运动（Brownian motion），具有相对较小的动量转移，大致与温度的量级相当，即$q^2 ∼ T^2$~\parencite{Dong:2019byy,Rapp:2009my}。由于能量传递按照$T/m_Q$进行参数压低，弹性相互作用占主导地位。
	重夸克的输运可以通过修正的朗之万方程来很好地描述~\parencite{Moore:2004tg,Cao:2013ita,Zhou:2016vwq,Dai:2018mhw,Wang:2019xey,Wang:2020ukj}：
  \begin{equation}
    	  	\vec{x}(t+\Delta t)=\vec{x}(t)+\frac{\vec{p}(t)}{E} \Delta t
    	  	\label{eq:Lange1}
   \end{equation}
  \begin{equation}
	  \vec{p}(t+\Delta t)=\vec{p}(t)-\Gamma \vec{p} \Delta t+\vec{\xi}(t)-\vec{p}_g
	  \label{eq:Lange2}
\end{equation}
	这两个方程分别描述了重夸克在穿越QGP介质时位置和动量随时间演化的信息。
  其中蒙特卡洛模拟过程中的最小时间间隔是 $\Delta t = 0.1 $fm；
  $\Gamma$即上一节介绍过的拖拽系数，决定了介质中重夸克与弹性碰撞能量损失的强度；
  $- \vec{p_g}$ 代表介质诱导胶子辐射而引起的动量修正；
  噪声项 $\vec{\xi}(t)$ 是引发扩散效应的随机力，描述了热化介质中准粒子对部分子的随机碰撞，且其概率分布服从高斯分布：
  \begin{equation}
  	W[\vec{\xi}(t)]=N \exp \left[\frac{\vec{\xi}(t)^2}{2 \kappa \Delta t}\right]
	\label{eq:Gau}
  \end{equation}
  该高斯分布的均值为0，方差为$\kappa \Delta t$。其中，$N$ 为归一化常数，$\kappa$为重夸克动量扩散系数。同时 $\vec{\xi}(t)$ 满足：
  \begin{equation}
  	\begin{aligned}
  		\left\langle\xi_i(t)\right\rangle & =0 \\
  		\left\langle\xi_i(t) \xi_j\left(t^{\prime}\right)\right\rangle & =\kappa \Delta t \delta_{i j}\left(t-t^{\prime}\right)
  	\end{aligned}
	\label{eq:xi}
  \end{equation}
  其中动量扩散系数$\kappa$与拖拽系数$\Gamma$存在耗散定理：
  \begin{equation}
    \kappa = 2\Gamma ET
  	\label{eq:kappa}
  \end{equation}
  所以Langevin方程的两个重要参数 $\kappa$ 和 $\Gamma$ 是伴随确定的，而其同时与空间扩散系数 $D_s$ 存在关系：
  \begin{equation}
	\kappa D_s = 2T^2
	\label{eq:Ds}
 \end{equation}
%张清
  其中 $D_s$ 可以由理论计算得到，在我们的研究中采用从 $D$ 介子产生的理论计算结果与实验数据比对得到的结果 $2\pi TD_s = 4$~\parencite{Rapp:2018qla}，这也满足格点量子色动力学 （Lattice QCD）理论计算给出的$2\pi TD_s \sim 3.7 - 7$ 的关系~\parencite{Banerjee:2011ra,Francis:2015daa,}。
  
对轻味部分子（轻夸克、胶子），我们采用了在“硬-热圈”（Hard-Thermal-Loop）~\parencite{Neufeld:2010xi,Huang:2013vaa}近似下的计算方法，用如下比率来描述轻味夸克和胶子的碰撞能量损失：
  \begin{equation}
  	\frac{d E}{d z}=-\frac{\alpha_s C_s \mu_D^2}{2} \ln \frac{\sqrt{E T}}{\mu_D}
	\label{eq:dE}
  \end{equation}
  其中强耦合常数 $\alpha_s = 0.3$，$z$是部分子沿传播方向的输运路径，德拜屏蔽质量（Debye screening mass）：
  \begin{equation}
  	\mu_D^2=4 \pi \alpha_s\left(1+\frac{n_f}{6}\right) T^2
  	\label{eq:Debye}
  \end{equation}
  $C_s$ 为特定色表示下的二次卡西米尔（Casimir）因子。
  % 对于一个李群，生成元是{J_1,J_2,...,J_i}，如果存在一个算符，使得[C,J_i]=0，则称C为Casimir算符。
%张清
    \begin{equation}
  	C_g = N_c
  	\label{eq:Cg}
  \end{equation}
  \begin{equation}
  	Cqg = \frac{N_c}{N_c^2-1}
	\label{eq:Cq}
  \end{equation}

以上即蒙特卡洛模拟的部分子在静态介质中弹性碰撞能量损失的情况。图~\ref{fig:elastic_EL}中为能量 $E_0 = 50 $ GeV 的胶子、轻夸克、粲夸克、底夸克在温度为 $T =400 $Mev 的静态介质中的弹性碰撞能量损失随时间$t$的变化。可以看到胶子的弹性碰撞能量损失要远高于轻夸克，同时轻夸克和重夸克虽然采用不同的模拟过程但呈现出相似的能量损失规律。蒙特卡洛模拟的结果与半解析计算的结果高度吻合。
  
  \hspace*{\fill}
  \begin{figure}[htbp]
  	\centering
  	\includegraphics[width=2.8in]{elastic_EL}
	\caption{比较蒙特卡洛模拟和半解析计算的结果，条件为 $400$ MeV温度的静态介质，以及具有初始能量 $E_0 = 50 $ GeV 的粒子的辐射能量损失~\parencite{Wang:2020qwe}}。
	\label{fig:elastic_EL}
  \end{figure}
  
  
  
%------------------------------- 2.7.2
\subsection{非弹性散射}
%辐射能量损失
在我们的框架中，采用了高扭度方法~\parencite{Guo:2000nz,Zhang:2003yn,Zhang:2004qm,Majumder:2009ge}来模拟喷注部分子在介质中诱导的胶子辐射，通过该方法可以获得QGP中高能部分子的胶子辐射谱：
  \begin{equation}
  	\frac{d N}{d x d k_{\perp}^2 d t}=\frac{2 \alpha_s C_s P(x) \hat{q}}{\pi k_{\perp}^4} \sin ^2\left(\frac{t-t_i}{2 \tau_f}\right)\left(\frac{k_{\perp}^2}{k_{\perp}^2+x^2 m^2}\right)^4
  	\label{eq:dN}
  \end{equation}
  其中 $x$ 和 $k_\perp$ 分别是辐射胶子的能量份额和横动量，强耦合常数 $\alpha_s$, Casimir 色因子 $C_s$ 取值如公式 \ref{eq:Cq}中表示，$M$ 是母部分子的质量。真空中劈裂函数 $P(x)$ ~\parencite{Deng:2009ncl}一般包含三个过程：
  \begin{equation*}
  	\begin{aligned}
  	& g \to g+g \\
  	& q \to q+g \\
  	& g \to q + \bar{q}
  	\end{aligned}
  	\label{eq:gqg}
  \end{equation*}
  其中胶子产生正反夸克对的过程（$g \to q + \bar{q}$）贡献很小，相对于另外两个贡献可以忽略，则有：
  \begin{equation}
    P_{q \rightarrow q g}(x)=\frac{(1-x)\left(1+(1-x)^2\right)}{x}
  	\label{eq:ggg}
  \end{equation}
  \begin{equation}
	P_{g \rightarrow g g}(x)=\frac{2\left(1-x+x^2\right)^3}{x(1-x)}
  	\label{eq:qqg}
  \end{equation}
%  真空中劈裂两数的表现形式图示化在图 3.6中，可以看到夸克劈裂函数随动量份额 $x$ 增大而下降，胶子劈裂函数关于 $x$ 呈现出对称分布。
% 图片略
  公式\ref{eq:dN}中 $\tau_f$ 是辐射胶子（daughter gluon）的形成的时间：
   \begin{equation}
  	  \tau_f = \frac{2Ex(1-x)}{k_\perp^2 + x^2m^2}
  	\label{eq:tau_f}
  \end{equation}
  两次胶子辐射时间间隔为 $t - t_i$。
  QGP中喷注输运参数 $\hat{q}$ 与温度 $T$ 存在关系 $\hat{q} \propto q_0(T/T_0)$~\parencite{Chen:2010te}，即
  \begin{equation}
  	\hat{q}(\tau, r)=\hat{q}_0 \frac{\rho^{Q G P}(\tau, r)}{\rho^{Q G P}\left(\tau_0, 0\right)} \frac{p^\mu u_\mu}{p^0}
  	\label{qhat}
  \end{equation}

  它决定了辐射能量的损失大小，取值范围可以由末态单举强子谱的理论计算与实验数据拟合得到。
	公式~\ref{eq:dN}中的最后一项是源自于重夸克“死角效应（“dead-cone” effect）”的压低因子~\parencite{ALICE:2019nuy}。
		重夸克的散射运动学与轻夸克的相比，重夸克通过介质时较不容易辐射能量~\parencite{Mustafa:1997pm}。这就导致了在发射角小于夸克质量与能量比值之间的角度内发射的胶子被压低（$\theta \lesssim m_Q/E$）~\parencite{Dokshitzer:2001zm,Armesto:2003jh}，形成了一个随着粒子能量减小而增长的“死锥”。这表明存在一种能量损失的层次结构$\Delta E_q>\Delta E_c>\Delta E_b$~\parencite{Xing:2019xae,Xing:2023ciw}，即与轻味强子相比，重味强子（及其衰变产物）的压低程度较小~\parencite{Buzzatti:2011vt,Kang:2014xsa}。
	
	
	
%------------------------------- 2.8
\section{本章小结}

	在本章讲述了微扰QCD理论的基本思想是利用微扰展开来计算物理过程的各阶贡献。在微扰QCD中，计算通常基于费曼图的形式，通过计算不同阶的修正项，以获得更精确的结果。
	其次，简要概述了高能核物理中用于识别和重构喷注的几种主要计算方法，它通过将能量沉积归类为单个喷注，并估计喷注的能量、动量和方向等物理量，并且也介绍了几种喷注重建算法的各自的特点和适用范围。
	
	然后阐述了蒙特卡洛事件产生器在高能核碰撞研究中被广泛用于模拟和生成实验中观测到的事件。这些事件产生器基于物理模型和随机数生成算法，可以生成各种粒子的产生、演化、相互作用和衰变等过程。
	接着介绍了一种常用的用于描述重离子碰撞的初始几何模型：Glauber模型。它基于核子的几何分布和碰撞概率，用于计算碰撞事件中核子的相互作用和空间分布。
	此外还介绍了在核物理和高能物理的研究中，部分子分布函数用于描述核子内部的夸克和胶子的携带动量分布。这些分布函数是实验数据和理论模型的结合产物，核修正后的部分子分布函数对于研究重离子碰撞、核物质的性质和高能物理实验中的过程具有非常重要的意义。
	最后，介绍了描述强子等粒子在强相互作用介质中输运行为的来两种模型：线性玻尔兹曼输运模型和SHELL。它们可以提供介质中粒子的演化和传播信息，为研究介质的输运性质、热力学性质和动力学行为提供理论框架。







\iffalse

  根据因子化(factorization)定理，可以把一个散射过程的微分截面分解成微扰可计算 的部分和非微扰可计算的部分。在微扰QCD中，非微扰可计算的部分通常是指分布函数和碎裂函数。分布函数(distribution function)是指在高动量的核子中发现携带特定动量的夸克的概率。	
  在深度非弹性散射实验中 $e+p \rightarrow e+X$ 利用因子化方案，整个散射过程截 面可以分解成微扰可计算的部分和非扰可计算的部分。
  部分子碎裂函数(Fragmentation function)，是指高能量的夸克或胶子碎裂成携带部分动量强子的概率。对于这些非微扰 的物理量，理论上我们可以通过微扰计算，得到它们随能量标度$Q^2$的演化方程。实验上 我们可以通过实验数据来提取它们的信息。

\section{脚注}
华中师范大学《附件4：关于修订毕业论文注释与参考文献著录格式的通知》提到

\begin{itemize}
  \item 文科术科的论文注释使用脚注
  \item 理工科的论文注释不使用脚注
\end{itemize}

\subsection{测试}
\subsection{测试}

\section{引用文中小节}\label{sec:ref}

如引用小节~\ref{sec:ref}

\section{引用参考文献}

这是一个参考文献引用的范例：“\parencite{邱泽奇建构与分化}提出……”。还可以引用多个文献：“\parencite{丁文祥rawtype,李晓东rawtype}提出……”。

英文文献 \parencite{feynman2011}

英文文献 \parencite[12]{feynman2011}

英文文献 \parencite[Thm1]{feynman2011}

英文文献 \parencite[12][Thm1]{feynman2011}

英文文献 \cite{feynman2011}

英文文献 \cite[12]{feynman2011}

英文文献 \cite[Thm1]{feynman2011}

英文文献 \cite[12][Thm1]{feynman2011}

\section{链接相关}

模板使用了 hyperref 包处理相关链接，使用 \verb|\href| 可以生成超链接，默认不显示链接颜色。如果需要输出网址，可以使用 \verb|\url| 命令，示例：\url{https://github.com}。
\begin{proof}
  \[
    x^2  \qedhere
  \]
\end{proof}

\fi